\chapter*{Abstract}

The acoustic impulse response is a central quantity in many applications, including the measurement of speech intelligibility and reverberation in rooms, of sound insulation between rooms, of sound attenuation outdoors, and of sound-speed profiles in the atmosphere. Measuring it has traditionally depended on complex wired equipment or high-latency Linux platforms, whose immobility and non-deterministic timing limit their use in distributed or outdoor scenarios. This thesis contributes to an affordable, wireless, synchronized system that acquires the responses to an arbitrary test signal between one source position and several receiver positions---which may be separated by a wall or by up to one kilometre---so that the impulse response can be recovered afterwards by post-processing on a host computer.

Each node combines a bare-metal STM32H7 microcontroller for precise audio capture with a CC1352P transceiver for wireless synchronization, keeping the time-sensitive signal path separate from the radio domain and using a custom event-based RF protocol to avoid operating-system delays. The contributions span the node subsystems: full-duplex audio was implemented and validated on a node, an analog front end was designed and built to bias and condition IEPE sensors, the sub-GHz link provides the wireless synchronization and range, and a custom power supply and node hardware were developed toward a single-board design.

The measurements show that the microcontroller audio path is deterministic to the sample, with a fixed, repeatable round-trip latency rather than the jitter that limited the earlier microprocessor-based system; because this latency is a known constant, it is simply accounted for during post-processing rather than appearing as timing uncertainty. The wireless link synchronizes the nodes to about one microsecond. Exponential-sine-sweep deconvolution recovers clean impulse responses with a high impulse-to-noise ratio and a frequency response flat across the audio band. The complete node was exercised acoustically, including reverberation-time and two-room sound-insulation measurements compared against a reference system; these confirm that the distributed nodes can act as an acoustic instrument, while identifying the custom power supply as the dominant limitation on the noise floor. The results show that a deterministic microcontroller node overcomes the timing limitations of earlier designs and supports autonomous, synchronized acoustic measurement across separated positions.